\section{Ingestion}
\label{sec:ingestion}

Ingestion turns an uploaded file into passages that carry their provenance.
It runs once per upload, in the order shown in \cref{lst:ingest}: the file is
read page by page, cleaned, split into labelled chunks, given the page numbers
printed on its pages, embedded and stored, and finally scanned for
cross-references (Part~II).

\begin{lstlisting}[float, language=Python, caption={The ingestion entry point
(\texttt{pipeline/pipeline.py}).}, label={lst:ingest}]
def ingest(doc_path: str) -> int:
    pages = load_document_pages(doc_path)
    chunks = with_printed_pages(chunk_document(pages),
                                load_printed_pages(doc_path))
    doc_name = os.path.basename(doc_path)

    embed_and_store_chunks(chunks, doc_name)
    build_graph(doc_name)

    return len(chunks)
\end{lstlisting}

\subsection{Reading documents}

A single loader dispatches on the file extension to a reader for PDF, DOCX or
plain text; the interface checks the extension before uploading, and the loader
rejects any other type. The loader returns a list of (page number, text) pairs,
so that page boundaries survive into the chunks. Only PDF files have real pages;
a Word or text file is returned as a single page.

PDF files are read with PyMuPDF, which reports every line of a page together
with its bounding box and font size. This geometry matters because of two kinds
of page furniture. Footnotes and running footers sit visually below the body
text, but once the layout is flattened into a text stream they land in the
middle of it: a footnote ends up spliced into the very clause that cites it,
where it corrupts both the passage and any reference found in it. The reader
therefore walks each page's lines from the bottom up and drops the trailing run
of lines that are either smaller than the body text or lie in the bottom 5\,\%
of the page, stopping at the first line of body text. The body size is taken to
be the font size that carries the most characters on the page. A recital marker
such as ``(148)'' is also set small, but it always has body text below it, so
the walk never reaches it. If PyMuPDF fails, the reader falls back to pypdf,
which extracts text without geometry.

\subsection{Cleaning}

Extracted text is normalised before it is parsed:

\begin{itemize}
  \item Unicode is normalised to NFC, and typographic dashes, quotes, bullets
  and ellipses are mapped to plain equivalents.
  \item Box-drawing and geometric symbols are removed, together with code
  points from the Unicode Private Use Area. Symbol fonts such as Wingdings
  place their bullets there, and extraction copies them as raw code points
  that render as empty boxes.
  \item Control characters and lines consisting only of a number are removed,
  runs of spaces are collapsed, and three or more blank lines become two.
  \item Text that lost its spaces during extraction---detected when fewer
  than 8\,\% of its characters are spaces---has letter runs of twelve or more
  characters split into words with a word-frequency model (wordninja).
\end{itemize}

Running headers and footers that survive the geometric filter are removed
across pages. The first and last four lines of every page are reduced to a
\emph{signature} in which every run of digits becomes \texttt{\#}, so that
``Page 12 of 144'' and ``Page 13 of 144'' share one signature. A signature
that appears on at least $\max(3, \lfloor P/2 \rfloor)$ of the document's
$P$ pages is page furniture and is removed everywhere before the pages are
joined. The threshold is deliberately conservative: a genuine line repeated
on half of all pages is rare, whereas a running header that survives turns up
in the middle of passages and in retrieval results.

\subsection{Structural units}

Densely cross-referenced documents announce their structure in their text:
headings such as ``Article 7'', ``Chapter III'' or ``Annex VIII'', recital
numbers such as ``(148)'', and numbered paragraphs and lettered points. The
structure parser finds these markers and the chunker turns them into chunk
boundaries and labels. This section describes the units down to the level of
a provision; the subdivisions inside a provision are the subject of
\cref{sec:structure}.

A heading is a unit word---\emph{Article, Section, Clause, Rule, Annex,
Schedule, Exhibit, Appendix, Chapter, Part} or \emph{Title}---followed by a
number in any of the forms the development corpora use: \texttt{58},
\texttt{3.13}, \texttt{3.7A}, \texttt{VII} or \texttt{B}. It must stand alone
on its line or be followed by a capitalised title, which rules out sentences
that merely begin with a reference. Article numbers must also form a sequence:
the first may be at most 3, and each later one must be greater than the one
before. A heading-like line that breaks the sequence---typically a reference
to another article that happens to begin a line---is rejected. Recitals are
numbers in parentheses alone on their line and are accepted only in the order
1, 2, 3, \ldots, which rejects footnote markers of the same shape.

Each unit has a depth (\cref{tab:depths}). The parser keeps the current path
from the outermost unit to the innermost, and a marker takes over its own
depth and closes everything deeper: a new article closes the previous
article's paragraph and point, but not its chapter. Every chunk is labelled
with the full path it falls under, such as \emph{chapter} III, \emph{article}
7, \emph{paragraph} 2, \emph{point} b.

\begin{table}[ht]
\centering
\small
\begin{tabular}{@{}cl@{}}
\toprule
Depth & Units \\
\midrule
0 & annex, schedule, appendix, exhibit \\
1 & part \\
2 & title \\
3 & chapter \\
4 & section \\
5 & article, recital, clause, rule \\
6 & paragraph \\
7 & point, subparagraph \\
8 & sub-point \\
\bottomrule
\end{tabular}
\caption{Structural depths. A marker closes every open unit at its own depth
or deeper.}
\label{tab:depths}
\end{table}

None of the unit words is specific to one legal system, and none is required.
A document in which fewer than three markers are found is treated as
unstructured and chunked plainly; retrieval then relies on similarity alone.

\subsection{Chunking}

The cleaned pages are joined into one string, and the offset at which each
page starts is recorded, so that any position in the text can be mapped back
to its page by binary search. The text between two consecutive markers is
split with a recursive character splitter into chunks of at most 500
characters with an overlap of 50, preferring to break at blank lines, then
line ends, then sentence ends, then spaces. Each chunk is stored with its
text, the pages on which it starts and ends, and the structural labels of its
path.

\paragraph{Trade-offs.} Small chunks make retrieval precise: one point of a
list is usually one chunk, so a citation points at the sentence that matters
rather than at a page of text. The cost is that a long provision spans many
chunks. This is acceptable because a named provision is never retrieved by
similarity: it is reassembled from its labels in reading order
(\cref{sec:retrieval}). Cutting at markers before splitting by size also means
that no chunk straddles two provisions, so every label a chunk carries is
exact.

\subsection{Printed page numbers}

The page number a reader sees is often not the PDF's own page index. Front
matter pushes the printed numbering behind: the page of TCPS~2 that prints
``48'' is the 56th page of the PDF. A reader who asks about ``page 48'' means
the printed number, and a citation should show it. The printed number sits in
the footer, which the reader has just removed, so it is recovered from the
page geometry instead.

For every page $p$, the numbers standing at the start or end of a line in the
top and bottom 8\,\% of the page are collected, including the first number of
forms such as ``117/144'' and ``L~176/5''. Each such number $n$ proposes an
offset $p - n$. Let $G_p$ be the set of offsets proposed on page $p$. The
document's offset is the one proposed by the most pages,
\begin{equation}
  o^{*} = \arg\max_{o} \bigl|\{\, p : o \in G_p \,\}\bigr|,
\end{equation}
and it is accepted only if at least $\max(3, P/2)$ of the $P$ pages agree on
it. Every page $p$ with $p - o^{*} \geq 1$ is then given the printed number
$p - o^{*}$, and every chunk receives the printed numbers of the pages on
which it starts and ends. For TCPS~2 the offset is 8; for the AI Act, whose
footers read ``117/144'', it is 0.

The agreement threshold is what keeps this general. Numbers in the margins
that are not page numbers---dates, article numbers, journal
references---propose offsets that vary from page to page and never
accumulate. A document
numbered per section (``3.8-2'') or not at all never reaches agreement and
keeps the PDF numbering, which is the correct fallback. The method assumes a
single dominant numbering run: front matter numbered in roman numerals simply
receives no printed number, and a document whose numbering restarts several
times is numbered by its longest run.

\subsection{Embedding and storage}

Chunks are embedded with all-MiniLM-L6-v2~\cite{wang2020minilm}, a
sentence-transformer~\cite{reimers2019sbert} that maps text to a
384-dimensional vector, and stored in a persistent ChromaDB collection that
uses cosine distance. Each chunk's identifier is
\texttt{\{document\}\_chunk\_\{i\}}, so reading order can be recovered by
sorting identifiers. Its metadata holds the source document, the start and end
pages (PDF and printed) and the structural labels. Plain numbers are stored as
integers and all other identifiers as text (\texttt{58} but \texttt{"3.7A"}),
because the store's equality filters are typed. Uploading a document again
first deletes its previous chunks, and writes are batched below the store's
maximum batch size, which large regulations exceed.

\paragraph{Trade-offs.} The embedder runs on the CPU. On a 4\,GB GPU every
megabyte of video memory taken by the embedder is a megabyte the language
model cannot use, and the model already has to offload part of its layers to
system memory. Embedding a large document on the CPU takes longer, but it
happens once per upload, and a query needs only one short text embedded.
ChromaDB was chosen over a bare vector index such as FAISS because it
evaluates metadata filters---equality, conjunction, disjunction and numeric
ranges---inside the store, which exact lookups (\cref{sec:retrieval}) depend
on. Its filters offer no substring or pattern matching, which is why the
parsing effort goes into producing exact labels at ingestion rather than into
matching loosely at query time.
